% =============================================================================
%  01_intro.tex —— §1 引言
% =============================================================================
\section{引言}\label{sec:intro}

我先为这个讲座谈的标题道个歉\cite{Ma}：标题是写给数学家的，讲稿本身却至少有一半
是讲给物理学家的。Kontsevich 的镜像对称纲领\cite{K} 与 D-膜（D-brane）的数学描述，
把三角范畴（triangulated category）带进了弦论与几何的主流；如今甚至有论文里的
[Ha] 指的是 [RD] 而不是 [AG]\yizhuRef{1}。这类范畴的名声一向吓人（多半是因为
参考文献大都是法文写的），而我的第二次报告还要用到它们，其他讲者想必也一样，
因此我想给这块漂亮的同调代数找一个动机。这个动机当然也是导出范畴的创立者
（主要是 Verdier，或者按这个圈子的传统叫法，Grothendieck 的学生 Verdier）
当初的动机；但同样当然，它没有被写下来\yizhuRef{2}。要了解理论的技术细节，
务实的数学家应当去读那本出色的 \cGM。

\begin{yizhu}[“$\mathrm{[Ha]}$ 指 [RD] 而不是 [AG]”]
这里说的是 Hartshorne 的两部书。\emph{Algebraic Geometry}（GTM 52，1977）
在文献里向来简称 [Ha] 或 [AG]；\emph{Residues and Duality}（LNM 20，1966）
简称 [RD]。作者的意思是：导出范畴成为常用语言之后，征引 Hartshorne 时，
更常见的反而是那本讲对偶与导范畴的小册子。
[来源：原文参考文献 HRD 条；GTM 52 书目信息属 A 级书目]
\end{yizhu}

\begin{yizhu}[Verdier 与“法文写的参考文献”]
Jean-Louis Verdier（1935--1989）在 1967 年的博士论文
\emph{Des catégories dérivées des catégories abéliennes} 中系统建立了导出范畴与
三角范畴，该文长期以油印本流传，1996 年才在 \emph{Astérisque} 239 刊行。
作者说“参考文献是法文写的”，指的就是这一系文献（Verdier 的论文，以及
SGA 4\textonehalf{} 中的相关章节）。
[来源：Verdier 博士论文，\emph{Astérisque} 239 (1996) 刊行；A 级]
\end{yizhu}

\begin{conceptbox}{导出范畴的一句话}
给定阿贝尔范畴 $\A$，其\term{导出范畴}{derived category} $D(\A)$ 是把
$\A$ 的复形范畴“翻转所有拟同构”得到的范畴：对象仍是复形，态射则允许沿拟同构
求逆。于是两个复形在 $D(\A)$ 中同构，当且仅当它们被一串拟同构连起来——
这正是“同调相同”的函子化版本，但比同调多留了信息。
[来源：原文 §\ref{sec:cat} 的定义 4.2；Gelfand--Manin 第 III 章]
\end{conceptbox}

导出范畴的主要想法很简单：\emph{与其余（上）同调打交道，不如与复形本身打交道}。
下面从代数几何取几个简单例子，说明为什么值得这样做；再从代数拓扑取几个例子，
说明这些想法与结构其实早已熟悉。（两门学科之间的联系是层，这里不展开。）
这也给了我们一套办法，把拓扑里的图形与技术搬进代数几何。

§\ref{sec:cx} 从代数几何与拓扑中的链复形讲起，说明它们为什么比（上）同调好用
\kanwuRef{1}。于是我们希望把空间、层之类的自然不变量取成复形，而不是它的同调；
要这么做，就得先规定哪些复形应当视为同构——这就引出\term{拟同构}
{quasi-isomorphism}的概念，即 §\ref{sec:qi} 的主题。建立一般（抽象）理论的地方，
范畴开始登场，见 §\ref{sec:cat}；由此得到锥、三角与三角范畴，见 §\ref{sec:cone}。
最后 §\ref{sec:df} 说明更经典的同调代数，特别是导出函子，如何纳入这个框架；
这使它们更透明，也让标准结果的证明与解释更容易。

\begin{kanwu}[引言中的节号错位]
原文作“Section 1 begins with chain complexes \ldots\ the subject of Section 2
\ldots\ in Section 3 \ldots\ in Section 4 \ldots\ Finally Section 5”，
即把六节按内容数成了 1--5，漏掉了引言本身。按实际的 \LaTeX{} 编号，
链复形是 §2、拟同构是 §3、范畴是 §4、锥与三角是 §5、导出函子是 §6。
译本正文改用实际节号（与原文其他地方的 \verb|\ref| 交叉引用一致），
详见附录 B 第 1 条。
\end{kanwu}

\textbf{致谢。}我从 Brian Conrad、Mikhail Khovanov 与 Paul Seidel 那里学到了
不少同调代数。感谢 CALF 代数几何讨论班的成员忍受这份讲稿初版的试讲\yizhuRef{3}，
也感谢 1999 年哈佛镜像对称冬季学校的组织者邀请我参加。

\begin{yizhu}[CALF]
CALF 是“junior \textbf{C}ambridge--\textbf{O}xford--\textbf{W}arwick”的缩写，
指剑桥、牛津、华威三校青年代数几何讨论班。原文把三个首字母加粗，正是为了点出
这个缩写。
[来源：原文正文的排版提示；A 级（原文自述）]
\end{yizhu}

\begin{keytakeaways}
  \item 全文只回答一个问题：为什么要把复形当作对象、把拟同构当作同构。
  \item 动机来自 Kontsevich 的同调镜像对称：D-膜的“导出范畴”不是技术装饰，
        而是物理上真的会区分同调看不出差别的复形。
  \item 读法有三条：拓扑（§2.1 → §5）、代数几何（§2.2 → §3 → §6）、范畴论（§4 → §5）。
  \item 原文引言的节号比实际少一（校勘 1），本书一律用实际节号。
\end{keytakeaways}
